\section{Rules for training reliably}
\label{sec:method}

The pattern in Section~\ref{sec:synthesis} suggests a straightforward
procedure, which we call CAS-Fuse. Four rules choose the settings that
trained reliably in the audit; two more keep several branches of the model
able to carry the degradation signal; and a short check after training
catches the seed dependence that no choice of architecture removes.

\subsection{How to set up the model}

\paragraph{Rule 1: window length by number of conditions} Window length
controls how many condition changes a training window averages over.
Single-condition subsets are insensitive to $w$; multi-condition subsets
were fragile in the audited mid-range (FD002 at $w\in\{30,60\}$, FD004 at
$w{=}40$). We use $w=40$ (FD001), $w=30$ (FD003) and the audited extremes
($w=20$ on FD002, $w=60$ on FD004).

\paragraph{Rule 2: depth by number of conditions} Changing depth does not
have the same effect at every condition count (Supplementary Table~S2):
FD002 $L\in\{2,4\}$ fail and $L\in\{1,3\}$ work; on FD004 $L{=}1$ fails
while $L\geq2$ works. We set $L=1$ for single-condition data and $L=3$ for
multi-condition data, outside the FD002 even-depth band that failed. The
resulting models have 1.08\,M and 1.47\,M parameters; all other
hyperparameters are shared, so $w$ and $L$ are the only choices that adapt.

\paragraph{Rule 3: keep more than one usable branch} The TCN front end,
difference augmentation and BiGRU stage stay in even where they add little
accuracy, because each is the branch the model depends on somewhere in the
range of condition structures: the TCN for short windows with many
conditions, attention over explicit rate signals for long windows with many
conditions, and the GRU as a residual backup. In every instrumented run it
was the convolutional branch that powered the recovery
(Section~\ref{sec:dynamics}). With several usable branches, failure needs
several independent things to go wrong at once.

\paragraph{Rule 4: keep the parts attention needs} Sinusoidal positional
encoding and sensor-channel dropout are retained: removing positional
encoding degraded mean RMSE by 16.9\% and raised the FD001 val--test gap from
17.5\% to 288.6\%; removing sensor dropout cost 3.1\% mean RMSE.

\subsection{What we removed, and why}
\label{sec:method_prune}

Three popular components were dropped on experimental grounds (full
removals in Supplementary Table~S6). Gradient reversal
\cite{ganin2016domain} looks like the answer to a condition cue and is
empirically the wrong one: it moves the failing settings around rather than
shrinking their number, helps by $\le$1.5 RMSE on single-condition data
where it is unnecessary, and injects condition-discriminating gradients into
the features most at risk of being picked up. Global and condition-aware
contrastive losses \cite{ragab2022contrastive} both made FD002 fail (43.46 /
43.49), since same-state pairs are rare under six interleaved regimes.
Removal A5 shows avg-only pooling slightly more accurate than the dual head
(16.43 versus 17.27 mean RMSE), but we keep the dual head because the
five-seed audit ran on this exact architecture; changing it would invalidate
the audit for a 4.9\% gain of unverified reliability. When reliability
matters, evidence about training should decide which components stay, not
only the accuracy difference.

\subsection{Checking a trained model}
\label{sec:method_audit}

No choice of settings removes seed dependence: the FD004 configuration from
Rules 1--2 failed on six of eight seeds while giving first-quartile accuracy
on the other two. The following checks are therefore part of the method,
and they are cheap.

\begin{enumerate}
  \item \emph{Train with $k\geq3$ seeds} (we use 5). This is the only step
        needing compute beyond standard practice.
  \item \emph{Check A: prediction std and Eq.~(\ref{eq:fingerprint}).}
        Flag if $\sigma_{\hat y}<0.01$ cycles, or if test RMSE matches the
        value the formula gives for a fixed output. One forward pass.
  \item \emph{Check B: attention entropy.} Flag if mean normalised entropy
        exceeds 0.99 (working 0.70--0.93; failed 1.0000). One hooked
        forward pass.
  \item \emph{Check C: a long flat stretch during training.} Flag a run as
        at risk when validation RMSE sits within $\varepsilon$ of
        $\operatorname{std}(y_{\mathrm{val}})$ with prediction std $<0.1$
        for $N$ consecutive epochs. The flat stretch is visible from epoch
        $\approx$5 and recovery happens by epoch 40, so at-risk runs can be
        stopped and started again with a new seed well before early stopping.
  \item \emph{Report mean $\pm$ std over working seeds together with the
        failure rate.} A configuration above a set threshold (we use 1/5) is
        rejected whatever its best-seed accuracy. On FD004, where every
        audited configuration depends on the seed, the procedure starts
        again with a new seed until Checks A--C pass, and reports the
        statistics of the runs that worked (19.53 $\pm$ 0.07 over two of
        eight).
\end{enumerate}

\subsection{What the rules do not promise}
\label{sec:method_tradeoff}

CAS-Fuse does not chase per-subset state-of-the-art accuracy; specialised
architectures reach lower single-run RMSE (Supplementary Table~S5). Within
the hybrid models used in applied RUL work, the claim is the best verified
reliability at competitive accuracy, together with an account of why the
rules work.
